\section{王朝不是在战场上建好的}

\subsection{甲子朝：牧野之后}

\eventanchor{WZ-1046-MUYE}{西周·牧野}

\begin{quote}
\centering\itshape “武王征商，唯甲子朝。”
\end{quote}

利簋把牧野留下来，只有这一句。没有阵势，没有喊杀，也没有后来史书里足以让人记住千年的细节；器物只记一个甲子日、一次“征商”、一次胜利。它像一枚冷硬的钉子，把周商之间最剧烈的断裂钉在时间里。约前1046年是今日常用的综合估算，不是能由一件器物独自裁定的年份；但“甲子朝”足以提醒人们，周人后来关于受命的宏大语言，立在一次真实而急遽的军事胜利上。

胜利也并不等于接管。商的都邑、祭祀、贵族、方国和道路仍在东方；认得商王、向商王纳贡或受其役使的人也仍在。周人则刚从渭水流域的政治网络中走出。打败一支军队是一回事，让比关中辽阔得多的地区承认新的册命、听从新的征发，又是另一回事。商末对外战争、王室与方国之间的离心、周方数代积累的联盟和资源，是牧野的结构性原因；牧野是这些裂缝在同一日断开的转折，并没有自动替胜者造出一个可以日常运转的王朝。

\phantomsection\label{fig:vol01:western-zhou:muye}
\begin{center}
\begin{tikzpicture}[x=.9cm,y=.9cm,font=\small]
  \fill[paper] (0,0) rectangle (10.2,5.5);
  \draw[blue!55,line width=1pt] (.4,4.5) .. controls (2.5,4.2) and (5.3,4.8) .. (9.8,4.2);
  \node[text=blue!70!black,above] at (5.0,4.45) {黄河};
  \draw[blue!45,line width=.75pt] (6.8,.6) .. controls (7.1,1.9) and (7.2,2.7) .. (7.1,3.8);
  \node[text=blue!70!black,right] at (7.2,2.3) {淇水方向};

  \fill[dynasty!23] (.8,1.0) -- (4.0,.9) -- (4.4,3.5) -- (1.3,3.7) -- cycle;
  \node[align=center] at (2.5,2.3) {周方\\自西而来};
  \fill[green!20] (7.4,1.5) -- (9.7,1.4) -- (9.8,3.6) -- (7.6,3.8) -- cycle;
  \node[align=center] at (8.6,2.65) {商都\\朝歌方向};
  \fill[timeline] (5.5,3.1) circle (.14);
  \node[above,align=center] at (5.5,3.2) {牧野\\战场区域};
  \draw[-{Stealth[length=2mm]},ultra thick,color=dynasty] (3.6,2.7) -- (5.35,3.05);
  \draw[-{Stealth[length=2mm]},thick,color=dynasty] (5.65,3.05) -- (7.25,2.8);
  \node[text=dynasty,below] at (4.6,2.65) {周方东进};
\end{tikzpicture}

\smallskip
\small\textit{图  牧野战事空间示意：突出周方由西向东进入商核心区、牧野与朝歌方向的关系；战场范围、行军路线与双方部署均为约略表达。}
\end{center}


\sourcebook{尚书·牧誓}把征伐说成天命转移，\sourcebook{史记·周本纪}又把它编进一条更完整的兴亡线。两种文本首先都在说明周何以有资格胜利。读者可以从中看见早周政治语言怎样把武力与秩序接在一起，却不能把它们当成已经摄下的战场实录。正是在这种限度内，牧野的意义才更清楚：它既是军事行动，也是新王室必须不断重述的合法性起点。

\begin{textwindow}[《尚书·牧誓》不是战地录音]
“牝鸡无晨。牝鸡之晨，惟家之索。”这句后来极有名的话见于《尚书·牧誓》。它把对商王朝的批判放进当时的宗法性别秩序中。今天读它，不是接受其中的价值判断，而是看周人怎样把军事行动说成恢复秩序的正当战争。
\end{textwindow}

\begin{sourcenote}[牧野能证实到哪里]
利簋是近时金文，为“甲子朝、征商”提供重要旁证；《尚书·牧誓》是传世政治文本，《史记》则是汉代对周初的回望。它们共同确认周取代商这一断裂，却不能合并成精确的战场日志。王年、甲子日与年代学模型相互校核后形成“约前1046年”的估算；兵数、阵势、逐句言辞和许多后世故事仍无从据此复原。
\end{sourcenote}

\subsection{监守旧都：征服者先借用被征服者}

\eventanchor{WZ-1046-THREE-GUARDS}{西周·三监}

灭商之后，周没有立刻把东方变成一片可以直接听命的王畿。传世叙事称武王使管叔、蔡叔、霍叔监守武庚与殷遗民。这一安排把宗室监督和旧商代理并置：前者提供周人的血缘与军事凭据，后者保留当地既有的人群、祭祀和统治关系。它不是一个从材料中可以画出职掌、辖区和俸禄的官制图，而是一种征服后不得不采取的权力拼接。

这种拼接带着内在的张力。若完全排斥商人，周必须在远离宗周的东方立即找到足够的人力与信息；若让武庚保有代理地位，又必须相信旧王朝的名望不会重新聚拢人群。管、蔡等宗室驻在其间，本可充作周王意志的延伸，亦可能成为不同判断、不同利益和不同不安的传递者。后来所谓“三监”的名称、人数和具体职权仍有文本层次问题，但旧商中心是周初最敏感的政治资源、周室曾以宗支与旧有力量同时处理它，这一基本处境不应被后出的整齐叙事遮蔽。

\eventanchor{WZ-1043-WUWANG-DEATH}{西周·武王崩}

武王不久去世，幼年的成王继位，原先暂时维系各方的安排突然失去中心。新王室面对的并不只是“谁来代政”的礼仪问题。册命、祭祀、军队调遣和对东方的承诺，究竟能否以幼主之名继续生效，决定着周人能否把牧野的胜利说成一个延续的王统。周公摄政因而不只是个人声望的故事；在可见的政治结构中，他必须让宗周、宗支和东方同时相信，王命没有因继承而中断。

王朝的日常面向在叙事缝隙间仍可辨认。周原甲骨记录祭祀、田猎与占问，\eventanchor{WZ-EARLY-ZHOUYUAN-ORACLES}{西周·周原甲骨}所留下的不是摄政会议的逐日纪录，也未必能逐片系到某一王年；它却表明早周的王室权威要靠反复的占问、祭祀和行动来维持。宏大的“受命”并非只在牧野或诰命中出现，也要在这种不断作出决定、确认关系的程序中才具有重量。

\subsection{摄政的风险：东方危机与宗支处分}

管叔、蔡叔、武庚以及东方势力随后被卷入反周危机。后来的史书很容易把他们排成“忠”与“叛”的队列；当时的人未必以同样清楚的道德语言理解自己的处境。旧商网络仍在，宗室之间对摄政的接受程度未必相同，幼主又不能亲自承担所有王命的可见动作。监守旧都原是稳定方案的一部分，在继承危机里却可能成为结盟、传递命令和彼此猜疑的渠道。

\eventanchor{WZ-1040-DONGZHENG}{西周·三监东征}

约前1042至前1036年的传统时段，较可靠的是相对次序而非一张按年排列的战事表：东方发生严重危机，周公一方发动东征，随后周重新安排东方关系。把这段过程称作“三监之乱、周公东征”，便利我们把握问题的规模，却不该暗示每次出兵的路线、兵力或参与者已经有同代军报可查。东征所要解决的直接问题，是摄政权威在东方能否生效；它的中期后果，则是周必须重新决定谁可以代表王室、谁可以保留旧有的地方资源。

\eventanchor{WZ-1040-WUGENG-SUPPRESSION}{西周·平定武庚}

传世材料与后出的叙事都把武庚及其盟附势力的失败置于这一过程之中。可以肯定的是，武庚失去了征服后被保留的代理位置，殷民的安置因而不再是附带问题，而成为周人必须正面处理的后果；不能肯定的，是每一战的地点、每一支力量的意图，或武庚在事件开端究竟说过什么、计划过什么。把沉默的材料补成戏剧性的密谋，只会使征服后的真正难题反而消失。

处分宗室同样不能只写成家族伦理的胜负。\eventanchor{WZ-1040-GUAN-SHU-PUNISHMENT}{西周·管叔被诛}，后世叙事记住了管叔被诛；\eventanchor{WZ-1040-CAI-SHU-EXILE}{西周·蔡叔被放}，又保留蔡叔被放逐、其后蔡仲续封的谱系。两种处置的审断过程、放逐地点和续封年份都不能由现有材料说死，但它们传出的政治讯号很明确：血缘不能自动豁免挑战王命的代价，宗支也并非只能一概清除。惩罚与再纳入并列，正是新王朝既要威慑、又不能耗尽可用亲族网络的选择。

\subsection{向东安置：一座城不能替代一张网络}

东征之后，周没有把东方仅仅当作战利品。传世文献把成周的营建、殷遗民的重新安置和东方封国的重编，连在平乱之后的政治图景中。这里最重要的不是把它想成一次完成的“迁徙”或一幅边界清晰的地图，而是看见周试图把城邑、受命家族、旧有人群、道路与祭祀关系编成新的秩序。成周作为东方支点，使宗周的王命能够更接近旧商腹地；受命者与被安置者之间如何分担劳役、守卫和祭祀责任，则仍须在地方关系中逐步落实。

\sourcebook{尚书}的《康诰》《酒诰》《梓材》保存的是以德、刑、祭祀和受命责任来说明治理的语言，不是一部现代意义的地方行政名册。大盂鼎等金文能够旁证早周处分、册命和政治关系的存在，也不能替我们列出东方每一处城邑、每一批人群的数量与去向。因此，“东部安置”可以说明周的国家能力开始依靠空间支点和代理网络扩展，不能被写成一项已经无缝完成的总规划。城邑把距离缩短，不能消除旧有记忆；封国分担统治成本，也保存了日后可以自行聚合的权力资源。

\begin{debate}[摄政、东征与安置能复原到什么程度]
《金縢》《大诰》《康诰》等篇的成文、传本与政治语境复杂；《史记》把它们组织为连贯的王朝叙事；金文与考古材料则提供不同尺度的旁证。三类材料的交集支持幼主继承、东方危机、严厉处置和重组关系这一相对次序，却不足以复原一份按日排列的摄政档案，也不足以断定管、蔡、武庚各自的内心动机。东征后的成周与安置，应当作为重建控制的过程来读，而不是一幕可被细节填满的凯旋场景。
\end{debate}

\subsection{归政：把非常之举交还给王权}

\eventanchor{WZ-1030S-ZHOUGONG-RETURN}{西周·周公归政}

在东征与东方重组告一段落之后，早期周政治传统把周公归政于成王作为摄政的收束。它解决的并非“贤臣是否愿意退让”这样单薄的问题，而是一个新王朝如何使非常时期的代理权重新回到王位：成年王需要亲自承接册命与祭祀权威，摄政者也必须让自己的权力不被理解为另一条王统。归政于是把危机中的必要授权，转换成王权可以延续的形式。

归政的具体年份仍有分歧，摄政持续多久也不能由《金縢》或《史记》直接算定。可是它作为政治记忆的力量是真实的：后人不断用周公的摄政与归政解释王权、宗支和辅政者应当如何相处。这种记忆本身不证明每一个情节，却说明周初的继承危机留下了一套可反复调用的合法性语言。牧野所开启的不是一条直线的胜利史，而是一连串必须被重新证明的命令关系。

\begin{causalchain}[从胜利到可持续的王命]
\eventref{WZ-1046-MUYE}{西周·牧野}解决了谁在战场上获胜；\eventref{WZ-1046-THREE-GUARDS}{西周·三监}把征服后的旧商中心置入宗室监督与地方代理的张力中；\eventref{WZ-1040-DONGZHENG}{西周·三监东征}则迫使周人以处置、安置和东方支点重建控制。至\eventref{WZ-1030S-ZHOUGONG-RETURN}{西周·周公归政}，非常摄政被交还给成王名下的王权。封国网络和成周支点提高了周室调动资源的能力，也留下地方家族、土地和军事关系能够积累的条件；这既是周初秩序的使能条件，也是它日后必须面对的制度代价。
\end{causalchain}
